\makeatletter
\def\input@path{{../styles/}{../../styles/}{../../../styles/}{../}{../../}{../../../}}
\makeatother
\documentclass{ee122_notes}
\input{macros.tex}
\input{preamble.tex}
\renewcommand{\releasedate}{April 20, 2026}
\begin{document}
\section*{EE 122/ME 141 Week 13, Lecture 1: Approximate Bode Plots}
\subsection*{Instructor: \instructor}
\subsection*{Date: \releasedate}
\section{Announcements}
\begin{itemize}
    \item Problem set \#10 due on April 22, 2026 
    \item Project milestone \#2 on designing a controller due on April 26 (Sunday) at 11:59 PM.
\end{itemize}
\section{Learning Objectives}
By the end of this lecture, you should be able to:
\begin{itemize}
    \item Understand how to sketch approximate Bode plots for transfer functions.
    \item Derive the ''rules of thumb'' for quickly sketching Bode plots.
    \item Appreciate the concept of loop shaping using Bode plots.
\end{itemize}

\section{Bode plots}
Last week, we introduced Bode plots as a graphical tool to analyze the frequency response of a system. It is a visualization that connects our understanding of the system to the real-world because it shows how the system responds to different frequencies of input signals. Bode plots consist of two separate plots: one for magnitude (in decibels) and one for phase (in degrees), both plotted against frequency on a logarithmic scale. The magnitude plot shows how much the system amplifies or attenuates signals at different frequencies, while the phase plot indicates how much the system shifts the phase of the input signal at those frequencies. Mathematically, for a transfer function \( H(s) \), the Bode plot is generated by evaluating \( H(j\omega) \) where \( \omega \) is the angular frequency. The magnitude is calculated as \( 20 \log_{10} |H(j\omega)| \) and the phase is calculated as \( \arg(H(j\omega)) \), that is the angle of the complex number \( H(j\omega) \). 
\subsection{Steps to compute Bode plots}
First, by substituting \( s = j\omega \) into the transfer function \( H(s) \), write \( H(j\omega) \), the frequency response of the system. Then, it is useful to write the expression of \( H(j\omega) \) as a complex number with a real part and an imaginary part. To do this, you can multiply the numerator and denominator by the complex conjugate of the denominator. This allows us to compute the magnitude and phase of the system at different frequencies. The magnitude is computed as \[ |H(j\omega)| = \sqrt{\text{Re}(H(j\omega))^2 + \text{Im}(H(j\omega))^2} \] and the phase is computed as \[ \angle (H(j\omega)) = \tan^{-1}\left(\frac{\text{Im}(H(j\omega))}{\text{Re}(H(j\omega))}\right) \]. 

But as you probably noticed, this process can be quite tedious. So, in this lecture, we discuss an approach to quickly obtain an approximate Bode plot without having to compute the exact magnitude and phase at every frequency. These are called asymptotic Bode plots as they are based on the asymptotic behavior of the system at low and high frequencies. The key idea is to break down the transfer function into its constituent parts, such as poles, zeros, and gain, and then apply simple rules to sketch the Bode plot based on these components.
\subsection{Sketching an approximate Bode plot}
Consider an example transfer function:
\[
G(s) = \frac{K}{s + a}
\]
where \( K \) is the gain and \( a \) is a positive constant. Let us apply the first step: write the frequency response by substituting \( s = j\omega \):
\[
G(j\omega) = \frac{K}{j\omega + a}
\]
Now, to compute the magnitude of the complex number, we remind ourselves of some properties that hold true for complex numbers:
\begin{align*}
|z_1 z_2| &= |z_1||z_2| \\
\left|\frac{z_1}{z_2}\right| &= \frac{|z_1|}{|z_2|} \\
\angle (z_1 z_2) &= \angle z_1 + \angle z_2 \\
\angle \left(\frac{z_1}{z_2}\right) &= \angle z_1 - \angle z_2
\end{align*}
Using these properties, we can write the magnitude and phase of \( G(j\omega) \) as follows:
\begin{align*}
|G(j\omega)| &= \frac{|K|}{|j\omega + a
|} \\
\angle G(j\omega) &= \angle K - \angle (j\omega +
a)
\end{align*}
which gives us 
\begin{align*}
|G(j\omega)| &= \frac{K}{\sqrt{\omega^2 + a^2}} \\
\angle G(j\omega) &= -\tan^{-1}\left(\frac{\omega}{a}\right)
\end{align*}
Let us now compute the gain in decibels. Recall that the decibel gain computation is simply \( 20 \log_{10} \) of the magnitude. This is a ``notational choice'' that has roots in sound and audio processing. Since it is a common measure for signals, Bode plots are also commonly plotted in decibels. Since the decibel gain computation involves a logarithm, let us remind ourselves of some properties of logarithms with the following popquiz.
\begin{popquiz}
For $a, b > 0$, write a simplified expression for each of the following:
\begin{enumerate}
  \item $\log (ab)$
  \item $\log \left(\frac{a}{b}\right)$
  \item $\log(a + b)$
  \item $\log (a^n)$
\end{enumerate}
\popqsplit 
\begin{enumerate}
  \item $\log (ab) = \log a + \log b$
  \item $\log \left(\frac{a}{b}\right) = \log a - \log b$
  \item $\log(a + b)$:  cannot be simplified further using logarithm properties! Don't confuse with the previous two properties.
  \item $\log (a^n) = n \log a$
\end{enumerate}
\end{popquiz}

Using the properties of logarithms, we can write the gain in decibels as follows:
\begin{align*}
20 \log_{10} |G(j\omega)| &= 20 \log_{10} \left(\frac{K}{\sqrt{\omega^2 + a^2}}\right) \\
&= 20 \log_{10} K - 20 \log_{10} \sqrt{\omega^2 + a^2} \\
&= 20 \log_{10} K - 10 \log_{10} (\omega^2 + a^2)\\
&= 20 \log_{10} K - 10 \log_{10} (a^2 (1 + \frac{\omega^2}{a^2})) \\
&= 20 \log_{10} K - 20 \log_{10} a - 10 \log_{10} (1 + \frac{\omega^2}{a^2})
\end{align*}

Let us denote the gain in decibels as \( |G|_{dB}\), so we have 
\[ 
|G|_{dB} = 20 \log_{10} K - 20 \log_{10} a - 10 \log_{10} \left(1 + \left(\frac{\omega}{a}\right)^2\right).
\] 
We can now analyze the behavior of the gain at low and high frequencies. Let us consider the following cases:
\paragraph{Case 1: Low frequencies ($\omega \ll a$)}
At low frequencies, we can approximate \( 1 + \left(\frac{\omega}{a}\right)^2 \approx 1 \) since \( \left(\frac{\omega}{a}\right)^2 \) is very small. Therefore, the gain in decibels can be approximated as follows:
\[
|G|_{dB} \approx 20 \log_{10} K - 20 \log_{10} a
\]
This means that at low frequencies, the gain is approximately constant and equal to \( 20 \log_{10} \left(\frac{K}{a}\right) \) dB.
\paragraph{Case 2: High frequencies ($\omega \gg a$)}
At high frequencies, we can approximate \( 1 + \left(\frac{\omega}{a}\right)^2 \approx \left(\frac{\omega}{a}\right)^2 \) since \( \left(\frac{\omega}{a}\right)^2 \) is very large. Therefore, the gain in decibels can be approximated as follows:
\begin{align*}
|G|_{dB} &\approx 20 \log_{10} K - 20 \log_{10} a - 10 \log_{10} \left(\frac{\omega}{a}\right)^2 \\
&= 20 \log_{10} K - 20 \log_{10} a - 20 \log_{10} \left(\frac{\omega}{a}\right) \\
&= 20 \log_{10} K - 20 \log_{10} a - 20 \log_{10} \omega + 20 \log_{10} a \\
&= 20 \log_{10} K - 20 \log_{10} \omega
\end{align*}
This means that at high frequencies, the gain decreases at a rate of 20 dB per decade (a decade is a tenfold increase in frequency).
\paragraph{Case 3: At the corner frequency ($\omega = a$)}
At the corner frequency, we can substitute \( \omega = a \) into the expression for the gain in decibels to find the gain at this frequency:
\begin{align*}
|G|_{dB} &= 20 \log_{10} K - 20 \log_{10} a - 10 \log_{10} \left(1 + 1\right) \\
&= 20 \log_{10} K - 20 \log_{10} a - 10 \log_{10} 2 \\
&= 20 \log_{10} K - 20 \log_{10} a - 3.01
\end{align*}
This means that at the corner frequency, the gain is approximately 3 dB less than the gain at low frequencies. This is called the 3 dB point, and it is a common reference point in Bode plots to indicate where the gain starts to decrease significantly (the bandwidth of the system).

So, we conclude the following important rule of thumb:
\begin{center}
  \textbf{Rule of thumb for sketching the effect of poles in Bode plots:} A pole contributes a slope of -20 dB/decade to the magnitude plot and a phase shift of -90 degrees at frequencies that are higher than the value at which the pole is located.
\end{center}

\subsection{Effect of zeros}
Let us first define zeros of transfer functions. For a general transfer function 
\[ 
G(s) = \frac{b_ms^m + b_{m-1}s^{m-1} + \ldots + b_1 s + b_0}{a_ns^n + a_{n-1}s^{n-1} + \ldots + a_1 s + a_0}
\]
the zeros are the values of \( s \) that make the numerator equal to zero, i.e., the roots of the numerator polynomial. With a simple example, let us derive a similar rule of thumb for sketching the effect of zeros in Bode plots. Consider the following transfer function:
\[
G(s) = s + b 
\]
where \( b \) is a positive constant. Substituting \( s = j\omega \), we can write the frequency response as follows:
\[
G(j\omega) = j\omega + b
\]
Using the properties of complex numbers, we can write the magnitude and phase of \( G(j\omega) \) as follows:
\begin{align*}
|G(j\omega)| &= \sqrt{\omega^2 + b^2} \\
\angle G(j\omega) &= \tan^{-1}\left(\frac{\omega}{b}\right)
\end{align*}
which gives us the gain in decibels as follows:
\begin{align*}
|G|_{dB} &= 20 \log_{10} |G(j\omega)| \\
&= 20 \log_{10} \sqrt{\omega^2 + b^2} \\
&= 10 \log_{10} (\omega^2 + b^2) \\
&= 10 \log_{10} (b^2 (1 + \frac{\omega^2}{b^2})) \\
&= 20 \log_{10} b + 10 \log_{10} (1 + \frac{\omega^2}{b^2})
\end{align*}
Now, we can analyze the behavior of the gain at low and high frequencies. Let us consider the same cases as before: low frequencies, high frequencies, and the corner frequency.

\begin{enumerate}
  \item For low frequencies, we can approximate \( 1 + \left(\frac{\omega}{b}\right)^2 \approx 1 \) since \( \left(\frac{\omega}{b}\right)^2 \) is very small. Therefore, the gain in decibels can be approximated as follows:
\[
|G|_{dB} \approx 20 \log_{10} b
\]
and the phase can be approximated as follows:
\[
\angle G(j\omega) \approx 0^\circ
\]
  \item For high frequencies, we can approximate \( 1 + \left(\frac{\omega}{b}\right)^2 \approx \left(\frac{\omega}{b}\right)^2 \) since \( \left(\frac{\omega}{b}\right)^2 \) is very large. Therefore, the gain in decibels can be approximated as follows:
\begin{align*}
|G|_{dB} &\approx 20 \log_{10} \omega
\end{align*}
and the phase can be approximated as follows:
\[
\angle G(j\omega) \approx 90^\circ
\]
  \item For the corner frequency, we can substitute \( \omega = b \) into the expression for the gain in decibels to find the gain at this frequency:
\begin{align*}
|G|_{dB} &= 20 \log_{10} b + 3.01
\end{align*}
and the phase can be computed as follows:
\[
\angle G(j\omega) = 45^\circ
\]
\end{enumerate}
So, we conclude the following important rule of thumb:
\begin{center}
  \textbf{Rule of thumb for sketching the effect of zeros in Bode plots:} A zero contributes a slope of +20 dB/decade to the magnitude plot and a phase shift of +90 degrees at frequencies that are higher than the value at which the zero is located.
\end{center}
\subsection{Loop shaping}
Now that you understand the contribution of poles and zeros, you can design controllers that shape the loop gain of the system to achieve desired performance specifications. Recall that the loop gain is the key transfer function that determines the closed-loop as well as the error transfer function. So, if you are able to get desired properties for the loop gain, you can achieve desired properties for the closed-loop and error transfer functions. For example, if you want to achieve a certain bandwidth for the closed-loop system, you can design a controller that adds a zero at the desired bandwidth frequency to increase the gain at that frequency and thus increase the bandwidth of the closed-loop system. See if you can apply this intuition in answering the following quiz question:
\begin{popquiz}
A plant has one pole. Your client wants a roll-off at high frequencies of -80dB/decade, that is, for each 10x increase in frequency the gain should fall by 80 dB. Select the statements that are TRUE about the controller.
\begin{enumerate}
  \item PID controller would do this well.

  \item The controller must have three poles.

  \item The controller must have two poles.

  \item PI controller would do this well.

  \item The controller can have zeros in addition to the poles but number of poles would increase then to compensate for the lead.

  \item The poles must all be with negative real parts.

  \item Some poles can have positive real part.

  \item The controller cannot have any zeros in addition to the poles.
\end{enumerate}
\popqsplit
Options (2), (5), and (6) are true. 
\end{popquiz}  
\begin{popquiz}
A plant $G(s)$ already has high-frequency slope -40 dB/dec. You want the loop transfer function $L(s)=C(s)G (s)$ to have high-frequency slope -20 dB/dec instead. Which of the following could help achieve this?
\begin{enumerate}
\item Add one zero in the compensator, active in that frequency range

\item Add two zeros in the compensator, active in that frequency range

\item Add one pole in the compensator, active in that frequency range

\item Make the controller contribute +20 dB/dec in that range
\end{enumerate}
\popqsplit 
Options (1) and (4) are correct.
\end{popquiz}
\section{Next lecture}
We will talk about loop shaping in more detail with a second-order system example.
\end{document}
